% Beamer slides: Fortran Types, Expressions, and Control Flow

\documentclass[10pt,aspectratio=169]{beamer}
\usepackage[utf8]{inputenc}
\usepackage{listings}
\usepackage{hyperref}
\usepackage{xcolor}
\usepackage{ragged2e}
\usepackage{graphicx}
\usepackage{tikz}
\usetikzlibrary{positioning,calc,decorations.pathreplacing}

\usetheme{Dresden}
%\usetheme{Bergen}

\title{Programming Techniques 2026--2027}
\author{Hannu Parviainen}
\institute{Universidad de la Laguna}
\date{\today}

\lstset{
  language=fortran,
  basicstyle=\ttfamily\scriptsize,
  columns=fullflexible,
  keepspaces=true,
  backgroundcolor=\color{white},
  keywordstyle=\color{blue},
  commentstyle=\color{gray},
  stringstyle=\color{red},
  showstringspaces=false,
  breaklines=true,
  frame=lines
}

\subtitle{Lecture 2: Fortran Types, Expressions, and Control Flow}

\begin{document}

% Title
\begin{frame}
  \titlepage
\end{frame}

\begin{frame}[fragile]{Example: Triangle Area (Module + Program)}
\begin{lstlisting}[language=Fortran]
module triangle_operations
  implicit none
contains
  function area(x, y, z)
    real :: area ! function type
    real, intent(in) :: x, y, z
    real :: theta, height
    theta = acos((x**2 + y**2 - z**2) / (2.0*x*y))
    height = x * sin(theta)
    area = 0.5 * y * height
  end function area
end module triangle_operations

program triangle
  use triangle_operations
  implicit none
  real :: a, b, c
  print *, 'Welcome, please enter the lengths of the 3 sides.'
  read *, a, b, c
  print *, "Triangle's area:", area(a, b, c)
end program triangle
\end{lstlisting}
\end{frame}


\begin{frame}{Coding Style (Recommendations)}
\begin{block}{Conventions}
\begin{itemize}
\item Always use \texttt{implicit none}.
\item Write keywords and intrinsics in lower case; use descriptive names (\texttt{n\_planets}, not \texttt{np}).
\item Use the symbolic relational operators (\texttt{<}, \texttt{==}, \texttt{/=}), not the legacy \texttt{.lt.}, \texttt{.eq.}, \texttt{.ne.}.
\item Use small, consistent indentation (2–4 spaces); indent blocks and units.
\item Include unit names in corresponding \texttt{end} statements.
\end{itemize}
\end{block}
\begin{block}{Note}
To fit code on slides, I have relaxed some conventions in teaching materials.
\end{block}
\end{frame}


\begin{frame}[fragile]{Source Form (Free Form)}
\begin{block}{Rules}
\begin{itemize}
\item Up to 132 characters per line.
\item \texttt{!} introduces a comment.
\item \texttt{\&} continues a logical line.
\item \texttt{;} separates multiple statements on one line.
\item Blanks are significant: they separate keywords and names (\texttt{integeri} is an error).
\end{itemize}
\end{block}
\begin{block}{Example}
\begin{lstlisting}[language=Fortran]
a = 5.2; m_star = 1.0        ! two statements on one line
p_years = sqrt(a**3 &        ! one statement on two lines
               / m_star)
\end{lstlisting}
\end{block}
\end{frame}

\section{Types and Declarations}

\begin{frame}{Intrinsic Types}
  \begin{block}{Default Intrinsic Types}
  \begin{itemize}
    \item \textbf{Character}
    \begin{itemize}
      \item \texttt{character :: spectral\_class}
      \item \texttt{character(len=12) :: star\_name}
    \end{itemize}
    \item \textbf{Logical}
    \begin{itemize}
      \item \texttt{logical :: is\_binary}
    \end{itemize}
    \item \textbf{Numeric}
    \begin{itemize}
      \item \texttt{integer :: n\_planets}
      \item \texttt{real :: mass}
      \item \texttt{complex :: visibility}
    \end{itemize}
  \end{itemize}
  \end{block}
\end{frame}

\begin{frame}{Literal Constants}
  \begin{block}{Examples}
  \begin{itemize}
    \item Integers: \texttt{4096}
    \item Reals: \texttt{1.0}, \texttt{1.989e30}, \texttt{-0.5}
    \item Logical: \texttt{.false.}, \texttt{.true.}
    \item Character: \texttt{"Barnard's Star"}, \texttt{'Barnard''s Star'} (escaped quote)
  \end{itemize}
  \end{block}
  \begin{block}{Notes}
  \begin{itemize}
    \item Only two logical values.
    \item Reals contain a decimal point (and may use exponential form).
    \item Character literals may be delimited by single or double quotes.
  \end{itemize}
  \end{block}
\end{frame}

\begin{frame}{Implicit Typing (Turn It Off)}
  \begin{block}{Rule}
  Without declarations, variables default to: \texttt{integer} if starting with I–N; otherwise \texttt{real}. This is error-prone (or, more plainly, stupid).
  \end{block}
  \begin{block}{Best Practice}
  Always put \texttt{implicit none} immediately after any \texttt{use} statements.
  \end{block}
\end{frame}

\begin{frame}[fragile]{Declarations (Numeric and Logical)}
  \begin{block}{General Syntax}
  \texttt{<type> [, <attributes>] :: <vars> [= <value>]}
  \end{block}
  \begin{block}{Examples}
\begin{lstlisting}[language=Fortran]
real :: mass
integer :: i, j
logical :: converged
real, dimension(3) :: position, velocity
integer :: n_steps = 1000
\end{lstlisting}
  \end{block}
  \begin{block}{Arrays}
  \texttt{dimension(3)} declares a three-element array (more on arrays in Lecture 3).
  \end{block}
\end{frame}

\begin{frame}[fragile]{Symbolic Constants (Parameters)}
  \begin{block}{Definition}
  Symbolic constants are declared with the \texttt{parameter} attribute.
  \end{block}
  \begin{block}{Examples}
\begin{lstlisting}[language=Fortran]
real, parameter :: pi = 3.14159
real, parameter :: G = 6.674e-11   ! m^3 kg^-1 s^-2
character(len=*), parameter :: host = 'Kepler-452', planet = "Kepler-452 b"
\end{lstlisting}
  \end{block}
  \begin{block}{When to Use}
  \begin{itemize}
    \item A variable should take only one value.
    \item To avoid “magic numbers”.
    \item For maintainability: values may change later.
  \end{itemize}
  \end{block}
\end{frame}

\begin{frame}[fragile]{Initialisation}
  \begin{block}{Rules}
  Initial values may use literals or other \texttt{parameter}s.
  \end{block}
  \begin{block}{Examples}
\begin{lstlisting}[language=Fortran]
real :: distance, luminosity = 3.828e26
integer :: n_stars = 5, n_steps = 100
character(len=1) :: band = 'V'
character(len=10) :: star = 'Vega'   ! right-padded with blanks
logical :: converged = .false., verbose = .true.
real, parameter :: pi = 3.141592, au = 1.496e11, year = 3.156e7
real :: v_earth = 2 * pi * au / year   ! Earth's orbital speed (m/s)
\end{lstlisting}
  \end{block}
  \begin{block}{Notes}
  A subset of intrinsics are allowed in initialisation expressions (e.g., \texttt{kind}, \texttt{len}, \texttt{huge}, \texttt{tiny}, \texttt{epsilon}, etc.).
  \end{block}
\end{frame}

\section{Expressions and Operators}

\begin{frame}{Expressions}
  \begin{block}{Types}
  Expressions can be numeric (e.g., \texttt{2*pi*r}), character (\texttt{star // ' b'}), or logical (\texttt{is\_binary .and. converged}).
  \end{block}
  \begin{block}{Use}
  Valid wherever an expression of the appropriate type is expected.
  \end{block}
\end{frame}

\begin{frame}[fragile]{Assignment}
  \begin{block}{Rule}
  The left-hand side is a variable; the right-hand side is an expression. The expression is evaluated first and then converted to the type of the variable. Numeric conversion happens silently.
  \end{block}
  \begin{block}{Examples}
\begin{lstlisting}[language=Fortran]
m_total = m1 + m2
x = r * cos(phi)                 ! cos uses radians
planet_name = star_name // ' b'
in_hz = (a >= a_inner .and. a <= a_outer)   ! in the habitable zone?
n_stars = 3.7                    ! integer n_stars: truncated to 3
\end{lstlisting}
  \end{block}
\end{frame}

\begin{frame}{Intrinsic Numeric Operations}
  \begin{block}{Operators}
  \begin{itemize}
    \item \texttt{**} exponentiation (right-associative): \texttt{10**2}
    \item \texttt{*}, \texttt{/} multiplication, division: \texttt{10*7/4}
    \item \texttt{+}, \texttt{-} unary/binary plus/minus: \texttt{10+7-4}, \texttt{-3}
  \end{itemize}
  \end{block}
  \begin{block}{Notes}
  Operators apply to literals, scalars, and arrays. With arrays they act element by element, and array operands must have the same shape.
  \end{block}
\end{frame}

\begin{frame}[fragile]{Relational Operators}
  \begin{block}{Syntax}
  \texttt{<}, \texttt{<=}, \texttt{==}, \texttt{/=}, \texttt{>=}, \texttt{>}\quad (legacy forms: \texttt{.lt.}, \texttt{.le.}, \texttt{.eq.}, \texttt{.ne.}, \texttt{.ge.}, \texttt{.gt.})\\
  Note that ``not equal'' is \texttt{/=}. Writing \texttt{!=} would start a comment.
  \end{block}
  \begin{block}{Example}
\begin{lstlisting}[language=Fortran]
is_giant = (radius > 10.0)
if (n_planets == 0) print *, 'No known planets'
\end{lstlisting}
  \end{block}
  \begin{block}{Reals}
  Avoid exact equality with reals; compare within a tolerance:
\begin{lstlisting}[language=Fortran]
real :: tol
tol = 1.0e-4
if (abs(m1 - m2) < tol) equal_mass = .true.
\end{lstlisting}
  \end{block}
\end{frame}

\begin{frame}{Logical Operators}
  \begin{block}{Operators}
  \begin{itemize}
    \item \texttt{.not.} — logical negation
    \item \texttt{.and.} — conjunction
    \item \texttt{.or.} — disjunction
    \item \texttt{.eqv.} — equivalence (same truth value)
    \item \texttt{.neqv.} — non-equivalence (different truth value)
  \end{itemize}
  \end{block}
  \begin{block}{Truth Table Highlights}
  \begin{itemize}
    \item \texttt{.not. .true.} is \texttt{.false.}
    \item \texttt{.true. .and. .false.} is \texttt{.false.}
    \item \texttt{.true. .or. .false.} is \texttt{.true.}
  \end{itemize}
  \end{block}
\end{frame}

\section{Control Flow}

\begin{frame}{Control Flow}
  \begin{block}{Overview}
    Control constructs allow the normal sequential order of execution to be changed. Fortran supports:
    \begin{itemize}
      \item Conditional execution: \texttt{if}, \texttt{if ... then ... else}, \texttt{select case}
      \item Loops: \texttt{do}, \texttt{do while}
      \item Loop control: \texttt{exit}, \texttt{cycle}
    \end{itemize}
  \end{block}
\end{frame}

\begin{frame}[fragile]{\texttt{if} Statement}
  \begin{block}{Syntax}
    \texttt{if (logical-expression) exec-stmt}
  \end{block}
  \begin{block}{Example}
\begin{lstlisting}[language=Fortran]
if (flux > flux_max) flux_max = flux
if (b_mag - v_mag > 1.0) is_red = .true.
if (n_obs > 0) mean_flux = total_flux / n_obs
\end{lstlisting}
  \end{block}
\end{frame}

\begin{frame}[fragile]{Block \texttt{if}}
  \begin{block}{Example}
\begin{lstlisting}[language=Fortran]
if (e < 1.0) then
  print *, "Elliptical orbit (bound)"
else if (e > 1.0) then
  print *, "Hyperbolic orbit (unbound)"
else
  print *, "Parabolic orbit"
end if
\end{lstlisting}
  \end{block}
  \begin{block}{Notes}
    Both \texttt{else} and \texttt{else if} are optional. Indentation improves readability.
  \end{block}
\end{frame}

\begin{frame}[fragile]{\texttt{select case}}
  \begin{block}{Example}
\begin{lstlisting}[language=Fortran]
select case (n_planets)
case (0)
  print *, "No known planets"
case (1)
  print *, "Single-planet system"
case (2:7)
  print *, "Multi-planet system"
case (8:)
  print *, "At least as many planets as the Solar System"
case default
  print *, "Invalid planet count"
end select
\end{lstlisting}
  \end{block}
  \begin{block}{Notes}
    Cleaner and often more efficient than nested \texttt{if}s.
  \end{block}
\end{frame}

\begin{frame}[fragile]{Indexed \texttt{do} Loops}
  \begin{block}{Example}
\begin{lstlisting}[language=Fortran]
do i = 1, 100, 1
  ...
end do
\end{lstlisting}
  \end{block}
  \begin{block}{Syntax}
    \texttt{do var = start, stop [, step]} with \texttt{step} defaulting to 1.
  \end{block}
\end{frame}

\begin{frame}[fragile]{Loop Variations}
  \begin{block}{Examples}
\begin{lstlisting}[language=Fortran]
do i = 1, 30, 2     ! 15 iterations
do j = 30, 1, -2    ! countdown
do k = 30, 1, 2     ! zero-trip (skipped)
do l = 1, 30        ! step default = 1
\end{lstlisting}
  \end{block}
\end{frame}

\begin{frame}[fragile]{Conditional \texttt{exit}}
  \begin{block}{Example}
\begin{lstlisting}[language=Fortran]
frame = 0
do
  frame = frame + 1
  if (frame > 100) exit
  print *, "Reducing frame", frame
end do
print *, "Loop finished. frame now equals", frame
\end{lstlisting}
  \end{block}
  \begin{block}{Explanation}
    A \texttt{do} without loop control repeats forever. \texttt{exit} jumps out of the current \texttt{do} loop.
  \end{block}
\end{frame}

\begin{frame}[fragile]{Conditional \texttt{cycle}}
  \begin{block}{Example}
\begin{lstlisting}[language=Fortran]
frame = 0
do
  frame = frame + 1
  if (frame >= 50 .and. frame <= 59) cycle   ! clouds: skip these frames
  if (frame > 100) exit
  print *, "Reducing frame", frame
end do
print *, "Loop finished. frame now equals", frame
\end{lstlisting}
  \end{block}
  \begin{block}{Explanation}
    \texttt{cycle} skips to the next iteration of the innermost loop.
  \end{block}
\end{frame}

\begin{frame}[fragile]{\texttt{do while} Loop}
  \begin{block}{Example}
\begin{lstlisting}[language=Fortran]
do while (.not. converged)
  ...
end do
\end{lstlisting}
  \end{block}
  \begin{block}{Equivalent}
\begin{lstlisting}[language=Fortran]
do
  if (converged) exit
  ...
end do
\end{lstlisting}
  \end{block}
\end{frame}

\begin{frame}[fragile]{Named and Nested Loops}
  \begin{block}{Example: counting bright pixels in an image}
\begin{lstlisting}[language=Fortran]
rows: do j = 1, ny
  cols: do i = 1, nx
    if (image(i, j) < 0.0) cycle rows          ! bad pixel: skip the rest of this row
    if (image(i, j) > saturation) exit rows    ! saturated pixel: stop searching
    if (image(i, j) < threshold) cycle         ! same as cycle cols
    n_bright = n_bright + 1
  end do cols
end do rows
\end{lstlisting}
  \end{block}
  \begin{block}{Notes}
    \texttt{exit} and \texttt{cycle} can name the loop they act on. Without a name, they act on the innermost loop.
  \end{block}
\end{frame}

\section{Wrap-up}

\begin{frame}{Takeaways}
  \begin{block}{Essentials}
  \begin{itemize}
    \item Use \texttt{implicit none}; declare all variables.
    \item Prefer symbolic constants for fixed values.
    \item Write keywords in lower case and use the symbolic operators (\texttt{<}, \texttt{==}, \texttt{/=}).
    \item Control flow constructs: \texttt{if}, \texttt{do}, \texttt{do while}, \texttt{select case}.
  \end{itemize}
  \end{block}
\end{frame}

\section{Exercises}

\begin{frame}{Exercises}
  \begin{block}{Instructions}
  \begin{itemize}
    \item Write each exercise as a separate program in \texttt{students/<initials>/sandbox/l02/}.
    \item Compile with checks on: \texttt{gfortran -Wall -fcheck=all -o prog prog.f90}
    \item The programs do not read input yet: set the input values with assignments at the start of the program. To test another case (including invalid values), edit them and recompile.
    \item Commit your programs and push them to your fork.
  \end{itemize}
  \end{block}
  \begin{block}{Overview}
  \begin{columns}[T]
    \column{0.48\textwidth}
    \begin{enumerate}
      \item Kepler's third law (arithmetic, conversion)
      \item Spectral classes (\texttt{select case})
      \item Habitable zone (logical operators)
    \end{enumerate}
    \column{0.48\textwidth}
    \begin{enumerate}
      \setcounter{enumi}{3}
      \item Wien's law (indexed \texttt{do})
      \item Kepler's equation (\texttt{do while})
    \end{enumerate}
  \end{columns}
  \end{block}
\end{frame}

\begin{frame}{Exercise 1: Kepler's Third Law}
  \begin{block}{Task}
  Set the semi-major axis $a$ (AU) of a planet and the mass $M$ ($M_\odot$) of its star in variables. Print the orbital period in years, in days, and as a whole number of days.
  \begin{itemize}
    \item Store the number of days in a year (365.25) as a \texttt{parameter}.
    \item Stop with an error message if either value is not positive (\texttt{stop 'message'}).
    \item Get the whole days by converting the period to an \texttt{integer}. Try both plain assignment and \texttt{int()}: what does \texttt{-Wall} say?
  \end{itemize}
  \end{block}
  \begin{block}{Physics}
  $P = \sqrt{a^3 / M}$, with $P$ in years, $a$ in AU, and $M$ in solar masses.
  \end{block}
  \begin{block}{Test}
  Jupiter: $a = 5.2$, $M = 1.0$ $\to$ $P \approx 11.86$ years $\approx 4331$ days.
  \end{block}
\end{frame}

\begin{frame}{Exercise 2: Spectral Classification}
  \begin{block}{Task}
  Set the effective temperature $T_\mathrm{eff}$ (K) of a star in an \texttt{integer} variable and print its spectral class using \texttt{select case}. Store the class letter in a \texttt{character(len=1)} variable and print it with \texttt{//}. Temperatures outside the sequence should give a message.
  \end{block}
  \begin{block}{Physics}
  \centering\setlength{\tabcolsep}{4pt}
  \begin{tabular}{ccccccc}
    O & B & A & F & G & K & M \\ \hline
    $\geq 30\,000$ & 10\,000--29\,999 & 7500--9999 & 6000--7499 & 5200--5999 & 3700--5199 & 2400--3699 \\
  \end{tabular}
  \end{block}
  \begin{block}{Test}
  Sun: 5772 K $\to$ G. \quad Vega: 9602 K $\to$ A. \quad 1500 K $\to$ outside the sequence.
  \end{block}
\end{frame}

\begin{frame}{Exercise 3: Habitable Zone}
  \begin{block}{Task}
  Set the luminosity $L$ ($L_\odot$) of a star, and the semi-major axis $a$ (AU) and eccentricity $e$ of a planet in variables. Report whether the planet is \emph{always}, \emph{partly}, or \emph{never} in the habitable zone (and if never, whether it is too hot or too cold).
  \begin{itemize}
    \item Use two \texttt{logical} variables: is the periapsis in the zone, and is the apoapsis in the zone?
    \item Hint: \texttt{.neqv.} is true when exactly one of its operands is true.
    \item Reject eccentricities outside $0 \leq e < 1$. What if the orbit crosses the whole zone?
  \end{itemize}
  \end{block}
  \begin{block}{Physics}
  Zone edges $r_\mathrm{in} = \sqrt{L/1.1}$ and $r_\mathrm{out} = \sqrt{L/0.53}$ AU; periapsis $a(1-e)$, apoapsis $a(1+e)$.
  \end{block}
  \begin{block}{Test ($L = 1$)}
  Earth (1.0, 0.017): always. \quad (1.3, 0.1): partly. \\ Venus (0.723, 0.007): too hot. \quad Mars (1.524, 0.093): too cold.
  \end{block}
\end{frame}

\begin{frame}{Exercise 4: Wien's Law Table}
  \begin{block}{Task}
  Print a table of the peak emission wavelength $\lambda_\mathrm{max}$ of a blackbody for $T = 3000, 6000, \ldots, 30\,000$ K, using an \texttt{integer} loop variable. Then print the table again, hottest first, using a negative step.
  \begin{itemize}
    \item Print the loop variable after each loop. Can you explain the values?
    \item What would you get if $b$ were the \texttt{integer} constant \texttt{2897772}?
  \end{itemize}
  \end{block}
  \begin{block}{Physics}
  $\lambda_\mathrm{max} = b / T$, with $b = 2.897772 \times 10^{6}$ nm K.
  \end{block}
  \begin{block}{Test}
  3000 K $\to$ 966 nm (infrared), \quad 6000 K $\to$ 483 nm (visible), \quad 30\,000 K $\to$ 96.6 nm (ultraviolet).
  \end{block}
\end{frame}

\begin{frame}{Exercise 5: Kepler's Equation}
  \begin{block}{Task}
  Set the mean anomaly $M$ (rad) and eccentricity $e$ of an orbit in variables, and solve Kepler's equation for the eccentric anomaly $E$ using Newton's method.
  \begin{itemize}
    \item Start from $E_0 = M$ and iterate in a \texttt{do while} loop until $|E_{n+1} - E_n| < 10^{-6}$ or 50 iterations have been done.
    \item Use a \texttt{logical} variable \texttt{converged}. Print $E$ and the number of iterations, or a warning if the iteration did not converge.
  \end{itemize}
  \end{block}
  \begin{block}{Physics}
  $E - e \sin E = M$, \qquad $E_{n+1} = E_n - \dfrac{E_n - e \sin E_n - M}{1 - e \cos E_n}$
  \end{block}
  \begin{block}{Test}
  $M = 1.0$, $e = 0.5$ $\to$ $E \approx 1.4987$. \quad $M = 1.0$, $e = 0.9$ $\to$ $E \approx 1.8621$.
  \end{block}
\end{frame}

\end{document}
